The partition and sharp-assignment conditions imply \(X_i\in\mathcal A_t\) if and only if \(T_i=t\). Condition on the scores, so every selected bandwidth is fixed. Each member of the triangular union belongs to a polynomial base class. For any boundary point \(x\), side \(t\), and radius \(r_0>0\), put \(s=\min\{h_0,r_0/2\}\). Positive polynomial lower mass gives
\[
 P_X\{X\in\mathcal A_t:d(X,x)<r_0\}\ge q_{t,P}(x,s)\ge c_ms^\kappa>0.
\]
Thus \(x\) belongs to each side support; Lipschitz continuity and the trace definition imply \(\theta_{t,P}(x)=g_{t,P}(x)\). Consequently, whenever the selected count \(N\) is positive,
\[
 \left|N^{-1}\sum_{i\in S}g_{t,P}(X_i)-\theta_{t,P}(x)\right|
 \le L\widehat h_{t,n}(x). \tag{81}
\]
If \(N=0\), the side fit is zero and the bounded-trace member property bounds its error by \(L\).

By \(\mathrm{lem:disk\text{-}pattern\text{-}sub\text{-}gaussian\text{-}control}\), at most
\[
 M_n=2|\mathcal H_n|\sum_{j=0}^3\binom nj
\]
nonempty averaging sets can occur. Under the exact conditional Gaussian property and conditional independence established in that lemma, its tail bound at \(z=\sqrt{2\log(2M_n/\gamma)}\) implies that, conditionally on the scores, with probability at least \(1-\gamma\), every nonempty candidate set satisfies
\[
 \left|N^{-1}\sum_{i\in S}\{Y_i-g_{T_i,P}(X_i)\}\right|
 \le\sigma\sqrt{\frac{2\log(2M_n/\gamma)}N}. \tag{82}
\]
Equations (81)--(82), together with the zero-count bound, show that each side error is at most \(B_{t,n}(x)\), simultaneously over \(x\) and \(t\). The triangle inequality proves conditional, and hence unconditional, simultaneous coverage. All counts and fits are measurable, the grids are finite, and the bandwidth rule has a deterministic tie break, so the band is data measurable. Since the same argument applies to every member of every component of the nonempty triangular union, taking the infimum over \(P\in\mathcal P_n(\Lambda_n)\) preserves the lower coverage bound \(1-\gamma\).

For the pervasive width assertion, put \(r=(\log n/n)^{1/(\kappa+2)}\) and choose \(h_*\in\mathcal H_n\) with \(r\le h_*<2r\). A maximal \(h_*/2\)-separated boundary net has at most \(2C_{\mathcal B}h_*^{-1}\) points. Every radius-\(h_*/2\) center-side count is binomial with mean at least \(nc_m(h_*/2)^\kappa\). For \(N\sim\operatorname{Bin}(n,q)\), Markov's inequality applied to \(2^{-N}\) gives \(\Pr(N\le nq/2)\le e^{-nq/8}\). A union bound and ball inclusion therefore give an event \(G_n\) such that
\[
 \inf_{x,t}N_{t,n}(x,h_*)\ge cn c_mh_*^\kappa,
 \qquad
 \sup_{P\in\mathcal P_\kappa^{\mathrm{perv}}}P(G_n^c)
 \le Ch_*^{-1}e^{-cnc_mh_*^\kappa}. \tag{83}
\]
Since \(nh_*^\kappa/\log(e+n/h_*)\to\infty\), the comparison bandwidth \(h_*\) is admissible everywhere on \(G_n\) for all sufficiently large \(n\). The minimizing property of \(\widehat h_{t,n}(x)\), together with \(nr^{\kappa+2}=\log n\), yields
\[
 L\widehat h_{t,n}(x)+
 \sigma\sqrt{\frac{\log(e+n/\widehat h_{t,n}(x))}{N_{t,n}(x,\widehat h_{t,n}(x))}}
 \le Cr \quad\text{on }G_n. \tag{84}
\]
Also \(\log(2M_n/\gamma)\le C\log n\), while \(\log(e+n/b)\ge\log n\) for \(b\le h_0<1\). Hence (84) implies \(B_{t,n}(x)\le Cr\) simultaneously on \(G_n\). On \(G_n^c\), each radius is at most \(L+\sigma\sqrt{C\log n}\): if \(N=0\) this is its definition, and if \(N>0\) it follows from \(N\ge1\) and \(\widehat h_{t,n}(x)\le h_0<1\). The exceptional contribution in (83) is \(o(r)\), because \(nh_*^\kappa\) grows polynomially in \(n\). Since the constructed interval has
\[
 W_n(x)=B_{0,n}(x)+B_{1,n}(x),
\]
the claimed uniform expected-sup width follows.

For the obstruction, fix \(x\in\mathcal B\) and let \(A_j\) be the event that \(\mathcal C_n(x)\) contains \(\tau_{P_j}(x)\). Honesty gives
\[
 P_j(A_j)\ge1-\gamma-\epsilon_n,
 \qquad j=0,1.
\]
By the definition of total variation for the full observed product experiments,
\[
 P_0(A_1)\ge P_1(A_1)-\operatorname{TV}(P_0^{\otimes n},P_1^{\otimes n}).
\]
Therefore the union bound gives
\[
 P_0(A_0\cap A_1)
 \ge1-2\gamma-2\epsilon_n-
 \operatorname{TV}(P_0^{\otimes n},P_1^{\otimes n}). \tag{85}
\]
On \(A_0\cap A_1\), the interval contains both target values, so its half-width is at least half their separation. Multiplying this implication by the positive part of (85) and taking \(P_0\)-expectations proves
\[
 \mathbb E_{P_0}W_n(x)\ge\frac{|\tau_{P_1}(x)-\tau_{P_0}(x)|}{2}
 \left[1-2\gamma-2\epsilon_n-
 \operatorname{TV}(P_0^{\otimes n},P_1^{\otimes n})\right]_+.
\]
This uses total variation of the full observed experiments. A score marginal is only a measurable projection and cannot supply the required upper bound on the full distance without an additional same-class outcome construction. Likewise, the global calibration has \(\log M_n\asymp\log n\), whereas \(\log\{eK_P(x,h)/\gamma\}\) may be bounded at an isolated thin site. The proof therefore establishes precisely the stated global band and full-law obstruction, but neither disclaimed sharper conclusion.